% id: 0542
% book: RPM 확률과 통계 · 06 통계적 추정
% section: 유형 05 표본평균의 확률; 미지수의 값 구하기
% figure: crop:fig-0542.png
% answer: $255$
% answer_source: 답지
% variant_level: 0 (원본 전사)
% 이 파일은 build.mjs 가 items.json 에서 만든다. 고칠 때는 items.json 을 고친다.
\smash{\raisebox{1.5mm}{\dmkichul{RPM 확률과 통계 0542번 · 92쪽}}}
\noindent\begin{minipage}[t]{0.58\linewidth}\vspace{0pt}\raggedright 정규분포 $\mathrm{N}(250{,}\mkern3mu\mathopen{}14^2)$을 따르는 모집단에서 크기가 $49$인 표본을 임의추출할 때, 표본평균 $\overline{X}$에 대하여 $\mathrm{P}(\overline{X}\ge k)\le 0.0062$를 만족시키는 실수 $k$의 최솟값을 오른쪽 표준정규분포표를 이용하여 구하시오.\end{minipage}\hfill\begin{minipage}[t]{0.4\linewidth}\vspace{0pt}\centering \setlength{\dmfigw}{106.7pt}\ifdim\dmfigw>\linewidth\setlength{\dmfigw}{\linewidth}\fi\includegraphics[width=\dmfigw]{figures/fig-0542.png}\end{minipage}\par
